\section{Theory extensions and recursion categories}
\label{sec:theory-extensions}

An inclusion of internal theories induces a fully faithful functor between
the Turing categories constructed in this section. Evaluation at a chosen
logical theorem compares a theory with its constant family in the
associated category. Under the deduction hypotheses of
Theorem~\ref{thm:associated-incompleteness}, a consistent theory is
essentially incomplete if and only if its constant family is.

\subsection{Categories associated with internal theories}

Let $\mathscr R$ be a recursion category, put
$\mathscr C=\Split(\mathscr R)$ and $\mathcal D=\Tot(\mathscr C)$,
and assume that $\mathcal D$ is regular with disjoint stable finite
coproducts. Assume also that every monomorphism $m:A\hookrightarrow X$
of $\mathcal D$ has a partial inverse $m^\dagger:X\rightharpoonup A$
with $m^\dagger m=1_A$ and $mm^\dagger=\overline{m^\dagger}$.
A monic span $X\xleftarrow{m}A\xrightarrow{f}Y$ in $\mathcal D$
represents $fm^\dagger$ in $\mathscr C$. Conversely, if
$h:X\rightharpoonup Y$ and $\bar h=mm^\dagger$, then $hm$ is total
and $h=(hm)m^\dagger$. Any two splittings of $\bar h$ give isomorphic
spans. For a second span $Y\xleftarrow{n}B\xrightarrow{g}Z$,
write $a:A\times_Y B\to A$ and $b:A\times_Y B\to B$ for the
pullback projections. The monomorphism $a$ represents the domain of
$n^\dagger f$. Indeed, a total $k:W\to A$ factors through $a$
if and only if $fk$ factors through $n$, equivalently if and only if
$n^\dagger f k$ is total. Hence $a a^\dagger=\overline{n^\dagger f}$.
Since $n^\dagger f a=b$, we have $n^\dagger f=b a^\dagger$ and
\[
 (g n^\dagger)(f m^\dagger)=(g b)(m a)^\dagger .
\]
The identity span represents the identity morphism, giving
$\mathscr C\simeq\Par(\mathcal D,\Mono)$.

For $T\in\mathcal D$, set
\[
 \mathscr C[T]=\Par(\mathcal D/T,\Mono).
\]
An object is a total morphism $p:X\to T$. A partial morphism from
$p:X\to T$ to $q:Y\to T$ is an isomorphism class of spans
$X\xleftarrow{m}A\xrightarrow{f}Y$ with $m$ monic and $pm=qf$.
Composition is by pullback. Its total category is $\mathcal D/T$.

\begin{lemma}[Turing structure in partial slices]
\label{lem:partial-slice-turing}
Under these assumptions, $\mathscr C$ and every $\mathscr C[T]$ are
Turing categories. If $Q$ is the original Turing object, with
$\tau:Q\times Q\rightharpoonup Q$, the Turing object in $\mathscr C[T]$
is $T\times Q\to T$. Under the product identification
$(T\times Q)\times_T(T\times Q)\cong T\times(Q\times Q)$, its
Turing morphism is
$1_T\times\tau:T\times(Q\times Q)\rightharpoonup T\times Q$ over $T$.
Let $\mathscr R[T]$ be the full restriction subcategory on the objects
$T\times X\to T$ for $X\in\mathscr R$. Then $\mathscr R[T]$ is a
recursion category and
\[
 \Split(\mathscr R[T])\simeq\mathscr C[T].
\]
\end{lemma}
\begin{proof}
The terminal object of $\mathcal D$ is restriction terminal in
$\mathscr C$. For $f:X\rightharpoonup Y$, split $\bar f=mr$.
The total morphism $fm$ satisfies $!_Yfm=!_Xm$, so composition with
$r$ gives $!_Yf=!_X\bar f$.

Every split object $(X,e)$ has a total embedding into $(X,1_X)$
and a partial retraction. An isomorphism $X\cong Q$ in $\mathscr R$
makes $(X,e)$ a retract of $Q$. The Turing morphism extends as in
Cockett--Hofstra~\cite[Section 3.5]{CockettHofstra}.
A morphism $(X,e)\times Q\rightharpoonup Q$ is represented in $\mathscr R$
by $f:X\times Q\rightharpoonup Q$ with $f(e\times1_Q)=f$.
Choose a total parameter $a:X\to Q$ with $f=\tau(a\times1_Q)$.
Then $f=\tau((ae)\times1_Q)$, and $ae:(X,e)\to Q$ is total since
$\overline{ae}=e$.
Thus $\mathscr C$ is Turing.

For $p:X\to T$, a total embedding $m_X:X\to Q$ gives a monomorphism
$m=\langle p,m_X\rangle:X\hookrightarrow T\times Q$. Its partial inverse
$m^\dagger:T\times Q\rightharpoonup X$ is a retraction over $T$.
The slice product of $p$ and $T\times Q\to T$ is
$X\times_T(T\times Q)\cong X\times Q$, with structure map $p\pi_1$.
Let $\varphi:(X\times Q,p\pi_1)\rightharpoonup(T\times Q,\pi_1)$
be a partial morphism over $T$. Its second component is
$f=\pi_2\varphi:X\times Q\rightharpoonup Q$ in $\mathscr C$.
Commutation over $T$ on the domain of definition gives
$\pi_1\varphi=p\pi_1\overline\varphi$.
Choose a total $a:X\to Q$ with $f=\tau(a\times1_Q)$. The total
parameter $\langle p,a\rangle:(X,p)\to(T\times Q,\pi_1)$ induces,
under the two slice product identifications, the morphism
\[
 \langle p\pi_1,\langle a\pi_1,\pi_2\rangle\rangle:
 X\times Q\longrightarrow T\times(Q\times Q).
\]
Its composite with $1_T\times\tau$ has second component $f$ and first
component $p\pi_1\bar\varphi$, so it equals $\varphi$.

An isomorphism $X\cong Y$ in $\mathscr R$ induces
$T\times X\cong T\times Y$ over $T$. The restriction product of two
such families is represented by $T\times(X\times Y)\to T$.
The zero span has domain the initial object of $\mathcal D/T$;
pullback stability of this initial object makes its product with any
partial map a zero span. Fullness supplies the parameter
$\langle p,a\rangle$ above whenever $p$ is a constant family.
For each $p:X\to T$, the graph embedding $m$ above satisfies
$m^\dagger m=1_X$ and $mm^\dagger=\overline{m^\dagger}$ over $T$.
Thus $m$ and $m^\dagger$ split the restriction idempotent
$\overline{m^\dagger}$ on $T\times Q\to T$ through $p$.
Since $\mathscr R[T]$ is full and $\mathscr C[T]$ is split,
the inclusion $\mathscr R[T]\hookrightarrow\mathscr C[T]$ induces
the claimed restriction equivalence.
\end{proof}

Fix a propositional or first-order deductive system in $\mathcal D$
as in Section~\ref{sec:internal-theories}, with sentence object $S$
and negation $\nu:S\to S$. In each slice $\mathcal D/P$, take the
syntax with constant family of sentences $P\times S\to P$ and the
logical rules with parameter $P$. Assume an initial derivation algebra
for every axiom subobject of this sentence family.
For a sentence theory $t:T\hookrightarrow S$, retain the sentence
family and theory
\[
 \widetilde S_T=(T\times S\xrightarrow{\pi_1}T),\qquad
 \widetilde t=1_T\times t:T\times T\hookrightarrow T\times S.
\]
The logical operations and rules retain the first coordinate.
Thus $\widetilde t$ is an internal theory in $\mathscr C[T]$.

\subsection{Functoriality under theory extension}

\begin{theorem}[Theory extensions and associated categories]
\label{thm:theory-extension}
Let $\mathscr R,\mathscr C,\mathcal D$ and the deductive system
satisfy the preceding assumptions. For internal theories
$t:T\hookrightarrow S$ and $u:U\hookrightarrow S$, let
$i:T\hookrightarrow U$ satisfy $t=ui$. Postcomposition defines a
fully faithful restriction functor
\[
 \Sigma_i:\mathscr C[T]\longrightarrow\mathscr C[U].
\]
Pullback defines $i^*:\mathscr C[U]\to\mathscr C[T]$, and the
underlying categories have adjunctions
\[
 \Sigma_i\dashv i^*\dashv\Sigma_i.
\]
If $i$ is proper, $\Sigma_i$ is not essentially surjective.
For composable theory inclusions $i,j$,
$\Sigma_j\Sigma_i=\Sigma_{ji}$.

The sentence and theorem comparison morphisms in $\mathcal D/U$
are $i\times1_S$ and $i\times i$. They give the commutative square
\[
\begin{tikzcd}
 \Sigma_i(T\times T\to T)
   \arrow[r,"i\times i"]\arrow[d,hook,"\Sigma_i\widetilde t"'] &
 (U\times U\to U)\arrow[d,hook,"\widetilde u"]\\
 \Sigma_i\widetilde S_T\arrow[r,"i\times1_S"'] & \widetilde S_U.
\end{tikzcd}
\]
Equivalently, $\widetilde S_T\cong i^*\widetilde S_U$ and
$\widetilde t\leq i^*\widetilde u$.
\end{theorem}
\begin{proof}
Postcomposition sends $p:X\to T$ to $ip:X\to U$ and retains each
span. Monicity of $i$ makes commutation over $U$ equivalent to
commutation over $T$ for spans between such objects. This proves full
faithfulness, restriction preservation, and the composition law.

For $p:X\to T$ and $q:Y\to U$, a span from $\Sigma_i p$ to $q$
factors its right leg uniquely through $T\times_UY$. This gives
\[
 \mathscr C[U](\Sigma_i p,q)\cong\mathscr C[T](p,i^*q).
\]
In the reverse direction let $Y\xleftarrow{m}A\xrightarrow{f}X$
represent $q\rightharpoonup\Sigma_i p$, so $qm=ipf$.
The pair $\langle pf,m\rangle:A\to T\times_UY$ is monic because its
composite with the projection to $Y$ is $m$. It is the left leg of a
span $i^*q\rightharpoonup p$. Conversely, composing a monic left leg
with $T\times_UY\hookrightarrow Y$ recovers a monic left leg over $U$.
These inverse assignments give
\[
 \mathscr C[T](i^*q,p)\cong\mathscr C[U](q,\Sigma_i p).
\]
Both assignments commute with pullback composition. They are natural,
and prove the adjunctions. Put $Y_T=T\times_U Y$, and write
$m:Y_T\hookrightarrow Y$ for the pullback of $i$. The unit of
$i^*\dashv\Sigma_i$ at $q$ is represented by the span
\[
 Y\xleftarrow{m}Y_T\xrightarrow{1_{Y_T}}Y_T
\]
over $U$. Thus the unit is $m^\dagger:q\rightharpoonup\Sigma_i i^*q$;
its restriction is $mm^\dagger$.

Every image object factors through $i$. If $1_U:U\to U$ were
isomorphic to one, the resulting factorization of $1_U$ through $i$
would make $i$ a split epimorphism and hence, by monicity, an
isomorphism. For proper $i$, the object $1_U$ therefore lies outside
the essential image of $\Sigma_i$.
The comparison identity
$(1_U\times u)(i\times i)=(i\times1_S)(1_T\times t)$
follows from $ui=t$. Pullback of the sentence family gives $T\times S$,
and its theorem family contains $T\times T$ by $1_T\times i$.
\end{proof}

Using the equivalences of Lemma~\ref{lem:partial-slice-turing},
the functor $\Sigma_i$ can be written
\[
 \Split(\mathscr R[T])\xrightarrow{\Sigma_i}\Split(\mathscr R[U]).
\]
For $p:X\to T$ and $q:Y\to T$, monicity of $i$ gives
$ipx=iqy$ if and only if $px=qy$ for every $x:W\to X$, $y:W\to Y$.
Thus $X\times_T Y\to X\times_U Y$ is an isomorphism preserving the
two projections, so $\Sigma_i$ preserves binary restriction products.
Its value at the restriction terminal object $1_T:T\to T$ is $i:T\to U$,
which is restriction terminal precisely when $i$ is an isomorphism.
For $q:Y\to U$, the unit $m^\dagger$ above is total precisely when
$m:T\times_UY\to Y$ is an isomorphism, that is, when $q$ factors through $i$.

\subsection{Essential incompleteness in the associated category}

Suppose a fixed logical theorem gives a total $\lambda:1\to S$
factoring through every theory under consideration. For
$u:U\hookrightarrow S$, write $\lambda_U:1\to U$ for its factorization.
Define
\[
 B_U:\mathscr C\to\mathscr C[U],\qquad
 B_U(X)=(U\times X\to U),\qquad B_U(f)=1_U\times f.
\]
Pullback along $\lambda_U$ sends $p:X\to U$ to
$E_U(p)=1\times_U X$ and pulls back both legs of a partial-map span.
Identifying $\mathcal D/1$ with $\mathcal D$ gives
$E_U:\mathscr C[U]\to\mathscr C$. The projections
$1\times_U(U\times X)\cong X$ give a natural isomorphism
$E_UB_U\cong1_{\mathscr C}$.
The isomorphism makes $B_U$ faithful. By construction $B_U$ preserves
restrictions, totality, and finite restriction products. It sends $\tau$ to $1_U\times\tau$, the Turing morphism
of Lemma~\ref{lem:partial-slice-turing}. It restricts to
$B_U:\mathscr R\to\mathscr R[U]$ on the original recursion category.

For a consistent theory $u$, let $\delta_u:S\rightharpoonup\Two$
have value $\mathbf t$ on the theorem subobject and $\mathbf f$ on
the refutation subobject $\nu^*u$. Their disjoint coproduct defines
this partial morphism, with restriction $u\vee\nu^*u$.
Subobjects are here identified with their restriction idempotents.
Put
\[
 \widetilde\delta_u=B_U(\delta_u):
 U\times S\rightharpoonup U\times\Two.
\]

\begin{theorem}[Incompleteness in an associated category]
\label{thm:associated-incompleteness}
Assume the categorical and syntax hypotheses of
Theorem~\ref{thm:theory-extension}. For every $P\in\mathcal D$ and
every axiom subobject $A\hookrightarrow P\times S$ in $\mathcal D/P$,
assume that the inference endofunctor for $A$ has an initial algebra.
Assume that a logical theorem $\lambda:1\to S$ factors through each
consistent internal sentence theory $u:U\hookrightarrow S$, and write
$\lambda_U:1\to U$ for its factorization. Then for every such $u$,
$u$ is essentially incomplete in $\mathscr C$ if and only if $\widetilde u=B_U(u)$ is
essentially incomplete in $\mathscr C[U]$. Moreover,
\[
\begin{aligned}
 u\text{ is incomplete}
 &\quad\Longleftrightarrow\quad
 \widetilde\delta_u\text{ is not total},\\
 (u,\nu^*u)\text{ is inseparable in }\mathscr C
 &\quad\Longleftrightarrow\quad
 \widetilde\delta_u\text{ has no total extension in }\mathscr C[U].
\end{aligned}
\]
If a consistent theory $t:T\hookrightarrow S$ admits a total reduction
of an inseparable pair to its theorem and refutation subobjects, every
consistent extension $u\geq t$ has an essentially incomplete theory
$\widetilde u$ in its associated category. The partial morphism $\widetilde\delta_u$ has no total extension.
\end{theorem}
\begin{proof}
Use the restrictions to total morphisms
\[
 B:\mathcal D\to\mathcal D/U,\qquad E:\mathcal D/U\to\mathcal D
\]
of $B_U$ and $E_U$, respectively. Both are changes of base, along
$!_U:U\to1$ and $\lambda_U:1\to U$, with $\mathcal D/1\cong\mathcal D$.
They preserve finite limits. Regularity makes images stable under pullback,
and extensivity makes finite coproducts stable under pullback. Thus
both functors preserve consistency and completeness.

Let $j:S\hookrightarrow\Phi$ include sentences among formulas.
For a closed sentence theory $v\hookrightarrow S$, let
$C(v)\hookrightarrow\Phi$ be the formula subobject generated by $v$,
so $j^*C(v)=v$. Applying $B$ to its rule operations equips $B(C(v))$
with the operations in the slice and includes the axioms $Bv$.
Initiality for $Bv$ therefore places every derivable formula in $B(C(v))$.
Restriction along $Bj$ and axiom introduction give
$\operatorname{Gen}(Bv)=Bv$.

In the other direction let $W\hookrightarrow BS$ be a closed sentence
theory in $\mathcal D/U$. Write $e:EB\cong1_{\mathcal D}$ for the
projection isomorphism. On the original sentence and formula objects put
\[
 W_0=(e_S^{-1})^*E(W),\qquad C_0=(e_\Phi^{-1})^*E(C(W)).
\]
Naturality gives $(EBj)e_S^{-1}=e_\Phi^{-1}j:S\to EB\Phi$.
Since $(Bj)^*C(W)=W$ and $E$ preserves pullbacks, $j^*C_0=W_0$.
The logical rules of the slice retain the parameter. Applying $E$ to
their operations on $C(W)$ and transporting along $e$ gives their
operations on $C_0$. This transport preserves premises and conclusions
by naturality of $e$ at the rule morphisms and the finite-product
comparisons. Thus $C_0$ is closed under the original rules and contains
the axioms $W_0$ as sentences. Initiality for $W_0$ and restriction along
$j$ give $\operatorname{Gen}(W_0)\leq W_0$. The reverse inclusion is
axiom introduction.

A consistent complete closed extension of $u$ therefore gives one of
$B_U(u)$ by $B_U$. Conversely, let $W\hookrightarrow U\times S$ be any consistent
complete closed extension of $B_U(u)$ in $\mathcal D/U$. Under
the identification $E B(S)\cong S$, its fiber is
\[
 W_0=\langle\lambda_U !_S,1_S\rangle^*W\hookrightarrow S ,
\]
where the pullback is taken in $\mathcal D$ along the displayed morphism
$S\to U\times S$. It is closed, consistent and complete by the
preceding argument. The inclusion $B(u)\leq W$ gives $u\leq W_0$.
This proves the essential-incompleteness equivalence.

The same retraction proves that $\delta_u$ is total exactly when
$B_U(\delta_u)$ is total. By definition, $\delta_u$ is total exactly
when $u$ is complete. The true and false fibers of a total extension of $\delta_u$ separate $u$ and
$\nu^*u$. Conversely, a complementary separator of these subobjects defines a
total extension of $\delta_u$. Either functor
carries a total extension to a total extension. The retraction therefore
proves the second equivalence.
Finally, Corollary~\ref{cor:domain-incompleteness} makes the theorem and
refutation subobjects of every consistent extension of $t$ inseparable.
Every consistent extension of $u$ also extends $t$, so $u$ is essentially
incomplete. The equivalence for constant families gives essential
incompleteness of $\widetilde u$, and inseparability of $u,\nu^*u$ excludes a
total extension of $\widetilde\delta_u$.
\end{proof}

In the numerical recursion category of partial recursive functions
(Section~\ref{subsec:numerical-realization}), let $A\subseteq\N$ be
infinite and computably enumerable.
Its partial identity $1_A$ splits through an object with carrier $A$.
To construct an isomorphism with $\N$, fix a program that halts on $x$
exactly when $x\in A$. Encode an input and a stage by a recursive
pairing and put
\[
 D=\{\langle x,s\rangle:\text{the program first halts on $x$ at stage $s$}\}.
\]
A simulation for $s$ steps decides membership in $D$. Each $x\in A$
occurs at exactly one such stage, so the first projection maps $D$
bijectively onto $A$. Since $D$ is infinite and recursive, its
increasing enumeration $d:\N\to D$ is total recursive. Thus
$e(n)=\pi_1(d(n))$ enumerates $A$ without repetitions. The inverse
search $r(x)=\min\{n:e(n)=x\}$ is partial recursive, with domain $A$.
As partial functions on $\N$ these maps satisfy
\[
     r e=1_{\N},\qquad e r=1_A .
\]
In the restriction splitting they give inverse total morphisms between
$\N$ and the object with carrier $A$. Postcomposition with these
morphisms identifies the slices over the two objects and hence their
categories of partial maps.

